% Pregunta 591 (2016-noviembre-OEP2016, n.º 36): la suma call larga + forward corto coincide con la put larga.
% Pagos a vencimiento con K = 3 (a escala):
%   suma: max(S_T-K,0) + K - S_T = max(K - S_T, 0) -> K - S_T si S_T <= K, 0 si S_T > K
% Autor: Víctor Gutiérrez Marcos
\documentclass[tikz,border=6pt]{standalone}
\input{tikz-tcee}
\begin{document}
\begin{tikzpicture}[tcee,x=0.95cm,y=0.95cm]
  \node[font=\small\bfseries,text=tceemorado,align=center] at (3,3.75)
    {put larga = call larga + forward corto};
  \draw[eje] (0,0) -- (6.3,0) node[below left] {$S_T$};
  \draw[eje] (0,-2.9) -- (0,3.3) node[left] {pago};
  \node[below left] at (0,0) {$0$};
  \node[left] at (0,3) {$K$};
  % Piezas, de fondo
  \draw[guia] (3,0) -- (5.8,2.8);
  \draw[guia] (3,0) -- (5.8,-2.8);
  \node[etiqueta,text=tceegris,font=\footnotesize,anchor=west] at (4.85,1.55) {$+(S_T-K)$};
  \node[etiqueta,text=tceegris,font=\footnotesize,anchor=west] at (4.85,-1.55) {$-(S_T-K)$};
  % Put comprada (ancha) y la suma encima
  \draw[tceemorado,line width=5pt,opacity=0.3] (0,3) -- (3,0) -- (5.9,0);
  \draw[curvanueva,desplazada] (0,3) -- (3,0) -- (5.9,0);
  \node[etiqueta,text=tceerojo,align=left,anchor=south west] at (1.55,1.6)
    {$\max(K-S_T,0)$};
  \node[etiqueta,text=tceerojo,font=\footnotesize,align=center,anchor=south] at (4.4,0.08)
    {se anulan: 0};
  \node[etiqueta,font=\footnotesize,align=left,anchor=north west] at (0.25,-1.3)
    {\textcolor{tceemorado}{banda: put larga}\\\textcolor{tceerojo}{discontinua: la suma}};
  \node[punto] at (3,0) {};
  \node[below=2pt] at (3,0) {$K$};
\end{tikzpicture}
\end{document}
